\section{Шаги сравнения и нормировка сертификата}
\label{app:explicit}
Если заданы $H_0,H_1$, рассмотрим
\[
 \alpha=\min\{1,G/[(H_0+H_1U)q^2]\},
\]
с полным шагом при нулевой кривизне. Аргумент спуска из
леммы~\ref{lem:local}, с учётом $U\ge\delta$, даёт
\[
 f_k-f_{k+1}\ge\frac12\min\{U_k,U_k^2/[D^2(H_0+H_1U_k)]\}.
\]
В фазе $b/2<U\le b$ продвижение не меньше
$\min\{b/4,b^2/(8D^2H_0+4D^2H_1b)\}$, поскольку правая часть
возрастает по $U>0$. Поэтому бюджет фазы допускает не более
$5+4H_1D^2+8H_0D^2/b$ обновлений. Суммирование по тем же
двоичным уровням доказывает
\begin{equation}\label{eq:explicitcount}
 N\le\ceil{(5+4H_1D^2)J+16H_0D^2/\varepsilon}.
\end{equation}
Для этого явного шага требуется $N+1$ вызовов каждого оракула.
Точная минимизация на отрезке даёт не меньшее продвижение, чем шаг
сравнения, поэтому наследует~\eqref{eq:explicitcount} без получения
$H_0,H_1$. Это не ограничивает внутренние вызовы функции или
градиента, необходимые для точного скалярного минимума. В расширенном
техническом отчёте также выведен гиперболический шаг того же порядка;
его элементарная функция не нужна для теоремы~\ref{thm:fw}.

Для полноты рассмотрим обратный поиск с нормировкой по сертификату:
он начинает с $\theta_{-1}=1$, пробует $\theta_{k-1}/2$, задаёт
$L=\theta U_k/D^2$ и удваивает $\theta$ при отказе квадратичной
проверки. Достаточный порог $T_k=D^2(H_0/U_k+H_1)$ не убывает.
Индукцией получаем
\[
 \theta_k\le\max\{2^{-k-1},2T_k\},\qquad
 L_k\le\max\{U_k/(2^{k+1}D^2),2(H_0+H_1U_k)\}.
\]
Действительно, приёмка первой пробы делит предыдущую границу пополам,
а после отказа принятая оценка меньше $2T_k$.
Значит, $L_k\le2H_0+(2H_1+1/(2D^2))U_k$. Применяя предыдущий
фазовый аргумент с этими эффективными константами, получаем
\[
 N\le\ceil{(7+8H_1D^2)J+32H_0D^2/\varepsilon}.
\]
Рекурсия $\theta_k=2^{r_k-1}\theta_{k-1}$ даёт
\[
 N_f\le1+2N+\ceil{\log_2\max\{1,2D^2(H_0/\varepsilon+H_1)\}}.
\]
В последней исходной точке перед обновлением $U>\varepsilon$,
что ограничивает $\theta_{N-1}$. Эта нормировка требует $D$ и
необязательна для двухфазного порядка. Теоремы~\ref{thm:fw}
и~\ref{thm:afw} относятся к обычному обратному поиску, а не к этому
отдельному варианту.

\section{Геометрические следствия и достигающие примеры}
\label{app:geometry}
Пусть $p=\inner{\nabla f(x)}{v-s}$ — парный зазор по полному
множеству атомов и текущему активному множеству. Выбранный зазор
FW/удаления не меньше $p/2$. Стандартное неравенство масштабирования
многогранника даёт для любого оптимального $x^*\ne x$
\[
 p\ge W\frac{\inner{\nabla f(x)}{x-x^*}}{\norm{x-x^*}}.
\]
См.~\cite[теорема 3]{lacostejaggi} и
\cite[леммы 3.4 и 3.6]{restartfw}. При квадратичном росте выберем
ближайший оптимум. Выпуклость и неравенство роста дают
$p\ge W\delta/\operatorname{dist}(x,X^*)
\ge W\sqrt{\gamma\delta/2}$. Поэтому
$g^2\ge\gamma W^2\delta/8$ и $\kappa=\gamma W^2/16$.
Монотонные итераты остаются в начальном множестве подуровня, где
предполагается рост. Единственность не требуется. При сильной
выпуклости
$\inner{\nabla f(x)}{x-x^*}\ge\delta+\mu\norm{x-x^*}^2/2$;
минимизация $\delta/r+\mu r/2$ по $r>0$ даёт
$p\ge W\sqrt{2\mu\delta}$ и $\kappa=\mu W^2/4$.
Эти геометрические следствия известны ранее.

\subsection{Равенство в огибающей зазора и нулевая базовая кривизна}
\label{app:bridge}
\paragraph{Квадратичное сравнение.}
При $C=D^2(H_0+H_1\delta)$ допустимый шаг
$t=\min\{1,G/C\}$ и лемма~\ref{lem:local} дают
$\delta\ge G^2/(2C)$ при $G\le C$ и
$\delta\ge G-C/2$ иначе. При $C=0$ используется полный шаг.
Обращение этого известного неравенства спуска даёт
\begin{equation}\label{eq:sharpbridge}
 G\le\widehat B(\delta):=
 \begin{cases}\sqrt{2C\delta},&2\delta\le C,\\
 \delta+C/2,&2\delta>C.\end{cases}
\end{equation}
Теорема~\ref{lem:bridge} сохраняет изменение кривизны вдоль отрезка
и усиливает это сравнение при $H_1>0$.

\paragraph{Достижимость при $H_1=0$.}
Пусть $Q=[-D/2,D/2]$, $x=D/2$ и $f(y)=h(y-z)^2/2$, $h>0$.
Для $z\in Q$ имеем $\delta=h(x-z)^2/2$ и
$G=hD(x-z)=\sqrt{2hD^2\delta}$, что достигает корневой ветви
\eqref{eq:sharpbridge} при $H_0=h,H_1=0$.
Для $z<-D/2$ условный оптимум равен $-D/2$,
$\delta=-hDz$ и $G=\delta+hD^2/2$, что достигает ветви полного
шага. Эти примеры также достигают предела точной огибающей при $H_1=0$.

\paragraph{Достижимость при $H_1>0$.}
Обозначим $k=\sqrt{H_1}$, $a=kD$, $c=H_0/H_1$ и возьмём
$Q=[0,D]$ с исходной точкой $x=0$. При $c>0$ и
$0<\delta\le c(\cosh a-1)$ положим
\[
 r=k^{-1}\operatorname{arcosh}(1+\delta/c),\qquad
 f(u)=c\big(\cosh(k(u-r))-1\big).
\]
Функция глобально выпукла, $f_Q^*=f(r)=0$, $f(0)=\delta$,
$f''=H_0+H_1f$, а градиент в исходной точке отрицателен.
Поэтому FW-атом равен $D$ и
$G=-Df'(0)=D\sqrt{2H_0\delta+H_1\delta^2}$.

При $\delta>c(\cosh a-1)$, включая все $\delta>0$ при $c=0$,
обозначим $E=E(\delta)$ и для $u\le D$ положим
\[
 f(u)=(\delta+c)\cosh(ku)-\frac{E}{a}\sinh(ku)-c.
\]
На $[0,D]$ выполняются $f(D)=0$ и
\[
 f'(D)=k\frac{c\cosh a-(\delta+c)}{\sinh a}<0.
\]
Тождество $f''=H_0+H_1f$ и эта производная в конце показывают,
что $f$ убывает и неотрицательна на отрезке: явно,
$f(u)+c=c\cosh(k(D-u))-[f'(D)/k]\sinh(k(D-u))$ при $u\le D$.
Она выпукла на всей левой полуоси. Для $u>D$ продолжим её как
\[
 f(u)=f'(D)(u-D)+\tfrac12H_0(u-D)^2.
\]
В точке $D$ совпадают значение, производная и вторая производная
обеих частей, поэтому продолжение глобально выпукло и принадлежит
$C^2$. Его условный оптимум равен $D$, $f(0)=\delta$ и
$G=-Df'(0)=E$. При $H_0=0$ получаем
$f(u)=\delta\sinh(k(D-u))/\sinh a$ для $u\le D$ с аффинным
продолжением за $D$. Эти конструкции доказывают оптимальность для
каждого набора параметров, а не только для фиксированной квадратичной
модели. Поскольку~\eqref{eq:sharpbridge} применимо к каждой
достигающей функции, они также доказывают
$E(\delta)\le\widehat B(\delta)$: новый порог никогда не меньше
квадратичного.

\paragraph{Точное глобальное следствие самоограничения.}
При другом предположении, когда условие гессиана выполняется на всём
$\R^d$ относительно конечного глобального инфимума $f^*$, имеем
\begin{equation}\label{eq:globalselfbound}
 \norm{\nabla f(x)}^2\le2H_0(f(x)-f^*)+H_1(f(x)-f^*)^2.
\end{equation}
Действительно, применим сравнение с положительным ядром на всём луче
$x-t\nabla f(x)/\norm{\nabla f(x)}$, $t>0$.
Неотрицательность $f-f^*$ ограничивает начальную направленную
производную величиной
\[
 \sqrt{H_1}\inf_{t>0}
 \frac{(f(x)-f^*+H_0/H_1)\cosh(\sqrt{H_1}t)-H_0/H_1}
 {\sinh(\sqrt{H_1}t)}
 =\sqrt{2H_0(f(x)-f^*)+H_1(f(x)-f^*)^2}.
\]
При $H_1=0$ используется квадратичный предел, а при $H_0=0<H_1$
— предел $t\to\infty$. Оценка достигается приведённой выше
глобальной cosh-функцией при $H_0,H_1>0$, квадратичной функцией
при $H_1=0<H_0$ и функцией $f(x)=A\exp(\sqrt{H_1}x)$,
$A>0$, при $H_0=0<H_1$; в последнем случае инфимум не достигается.
Это даёт точные константы глобальной модели.
Оценка~\cite[приложение G, лемма 6]{warmup} имеет большие
коэффициенты $9H_0/4$ и $27H_1/4$; точность там не заявляется.
Для аффинного ограничения второго порядка
\cite[следствие 1]{caoselfbound} вместо этого даёт
$\norm{\nabla f}^2\le2(H_0+H_1\delta)^2
\log(1+H_1\delta/H_0)/H_1$ при $H_0,H_1>0$.
Точное полиномиальное выражение~\eqref{eq:globalselfbound} устраняет
эту логарифмическую добавку. Сам принцип самоограничения известен
ранее~\cite[лемма 11]{kale}; здесь сравниваются лишь приведённые
константы и экстремизаторы. Одна лишь условная модель не
влечёт~\eqref{eq:globalselfbound}: непостоянная аффинная функция на
компактном отрезке имеет $H_0=H_1=0$ и ненулевой градиент.

\paragraph{Полиномиальный пример с нулевой базовой кривизной.}
Глобально выпуклый $C^2$-пример $f(x)=t+(\max\{t,0\})^3$,
$t=1/2-x\in[0,1]$ на $Q$, удовлетворяет
$f''(x)=6t\le6f(x)$, поэтому $H_0=0,H_1=6$.
Он не является сильно выпуклым. В неоптимальной внутренней точке
оба конца-атома активны, и выбранный зазор равен
$(1+3t^2)\max\{t,1-t\}\ge1/2$; неоптимальный конец также
удовлетворяет этому неравенству. Поскольку $\delta\le2$,
условие~\eqref{eq:geometry} выполнено с $\kappa=1/16$.
При нулевой ошибке геометрическое неравенство очевидно.
Это показывает невакуумность расширения на $H_0=0$.

\section{Контроли нижних границ и детерминированные ошибки оракулов}
\label{app:monitors}
Обозначим $g_j=\nabla f(x_j)$ и $a_j=f(x_j)-\inner{g_j}{x_j}$.
Для завершённого блока $B$ из $K$ последовательных оплаченных
исходных точек
\[
 b_B=\bar a_B+\min_{v\in Q}\inner{\bar g_B}v\le f_Q^*.
\]
Выпуклость даёт $f(v)\ge a_j+\inner{g_j}v$ для каждой точки;
усреднение и минимизация доказывают корректность.
Непересекающиеся блоки начинаются с исходной точки ноль;
незавершённый блок не вызывает оракул. Новая нижняя граница —
максимум $\ell_k$ и всех завершённых $b_B$.
На том же монотонном пути её сертификат не возрастает и не больше
$U_k$, поэтому первое пересечение порога происходит не позже по
обновлениям и вызовам функции или градиента. В исходной точке $k$
выполнено ровно $\lfloor(k+1)/K\rfloor$ дополнительных полных LMO.
Суммарная стоимость LMO не обязана уменьшаться. Мониторы не влияют
ни на линейный поиск, ни на начальное $M$; каждая указанная стоимость
относится к одному отдельному монитору без остальных диагностических
мониторов. Доступная по бюджету исходная точка оплачивает всю
запланированную проверку.

Эти равномерные блочные границы отличаются от двойственной величины
Фенхеля с усреднённым градиентом из~\cite[теорема 3]{pedregosa}.
Когда существует глобальная сопряжённая функция, равенство Фенхеля
и неравенство Йенсена дают
\[
 b_B=-\overline{f^*(g_j)}-\sigma_Q(-\bar g_B)
     \le-f^*(\bar g_B)-\sigma_Q(-\bar g_B).
\]
Здесь $f^*$ обозначает сопряжённую функцию, а не скалярный оптимум
$f_Q^*$. Перенос теоремы или весов усреднения той работы не
заявляется. При заданной корректной константе сильной выпуклости
$\mu$ оценка $f(x_j)-\norm{g_j}^2/(2\mu)\le f_Q^*$ следует из
минимизации квадратичной нижней модели на $\R^d$~\cite[приложение B]{karimi}.
Сохранение этих границ не требует LMO, но использует дополнительную
информацию. Нативные проверки остановки на одном синтетическом
внутреннем экземпляре проверяют все 20 пар «метод/профиль» по
конструкции достижимого префикса и инструментируют фактические
вызовы оракулов. Это не 2160 дополнительных нативных запусков решателя.

\subsection{Корректность сертификата при ограниченных детерминированных ошибках}
\label{app:inexact}
Если $|\widehat f-f(x)|\le\nu$,
$\norm{\widehat g-\nabla f(x)}\le\rho$ и
$\inner{\widehat g}s\le\min_{v\in Q}\inner{\widehat g}v+\eta$
с известными границами, то
\[
 \ell^{\rm safe}=\widehat f+\inner{\widehat g}{s-x}
                      -\nu-\rho D-\eta\le f_Q^*.
\]
Для каждого $v\in Q$ выпуклость даёт
$f(v)\ge\widehat f-\nu+\inner{\widehat g}{v-x}-\rho D
\ge\ell^{\rm safe}$; минимизируем по $v$.
Максимум таких нижних границ остаётся корректным, а
$\widehat f+\nu-\max_j\ell_j^{\rm safe}$ является верхним
сертификатом. Это утверждение о корректности, а не теорема сходимости
для шумных направлений. Эмпирическая дисперсия не является
детерминированной границей нормы.

\section{Дополнительные подробности протокола и контроли}
\label{app:protocol}
\paragraph{Признаки.}
Seoul использует восемь численных погодных признаков: температуру,
влажность, скорость ветра, видимость, точку росы, солнечную радиацию,
дождь и снег; синус/косинус часа и дня недели; индикаторы праздника
и функционирования. Metro использует температуру, дождь, снег,
облачность, синус/косинус часа, дня недели и дня года с периодом
365.25, а также индикатор праздника. Остальные текстовые категории
опущены. Постоянные на обучении признаки удаляются до преобразования
$\tanh$; кодировки источников и имена признаков записаны в манифестах.
Эти детерминированные решения принимаются до сравнения решателей.

\paragraph{Исходный вариант сравнения.}
Для $x\ne y$ оценка AC-FW равна
\[
 \ell(x,y)=\frac{2|f(y)-f(x)-\inner{\nabla f(x)}{y-x}|}{\norm{y-x}^2},
\]
а при $x=y$ — нулю. При обновлении берётся максимум этой пробной
оценки и $r_kL_k$, где $r_k=1-1/((k+1)\log^2(k+3))$.
Проба принимается лишь при строгом уменьшении функции.
После $N$ внешних проб оплаченные счётчики равны
$N_f=N_{\rm LMO}=N+2$, $N_g=N_{\rm accepted}+2$.
Отказы расходуют лимит. Все основные старты счётных данных — атомы,
как в цитируемой инициализации. Тесты со скрытыми параметрами на
обоих наборах данных инструментируют предварительную обработку,
отказы и вызовы оракулов независимо от записанных счётчиков.

\paragraph{Заданная геометрия.}
Для контроля инициализации положим
\[
 \phi(x)=m^{-1}\sum_i[\cosh(a_i^\top(x-x^*))-1]
                         +\mu\norm{x-x^*}^2/2,
 \quad f(x)=\phi(x)-\lambda v^\top(x-x^*),
\]
где $\norm{x^*}_1=R$, $v\in\partial\norm{x^*}_1$ и
$\lambda=0.15$. На $Q$ имеем $f\ge\phi\ge0=f_Q^*$ и
$\nabla f(x^*)=-\lambda v\ne0$. Корректные константы:
$H_0=\mu+\lambda_{\max}(A^\top A/m)$ и
$H_1=\max_i\norm{a_i}^2$.
Семейства active и active-scaled используют $d=12,m=10$,
$R=1.5,\mu=0.08$ с координатами $x^*$, равными $(0.6,-0.5,0.4)$
на позициях $(2,6,10)$; гауссовские элементы строк имеют стандартное
отклонение $0.22$ или $0.9$. Многомерное семейство имеет $d=60,m=40$
и стандартное отклонение $0.22$. У каждого десять согласованных seed.
Координаты нормали вне носителя равномерны на $[-0.6,0.6]$.

При $w=\operatorname{sign}(x^*)$, равном нулю вне носителя,
запас оптимальной грани $\sigma(x)=R-w^\top x$ удовлетворяет
$\sigma(x+\alpha(s-x))=(1-\alpha)\sigma(x)+\alpha\sigma(s)$.
Поэтому шаги строго меньше единицы к атомам грани сохраняют
положительный запас при старте вне грани. Полный начальный переход
к вершине может его устранить. Кроме того, на этих заданных
примерах $f(x)\ge0.4\lambda\sigma(x)$.
Это известное препятствие грани для FW объясняет необходимость
контроля инициализации, а не новый механизм ускорения.

\paragraph{Исследование единиц целевой функции.}
Отдельный архив радиусов/масштабов содержит 360 запусков на
логистическом контроле WDBC и пуассоновской функции UCI Bike
Sharing~\cite{wdbc,bike}: три радиуса, пять seed, три положительных
масштаба и четыре BT-метода. Точность масштабируется вместе с функцией.
60 парных сравнений масштабов на каждый метод с инициализацией по
сертификату сравнивают конечные счётчики и сертификаты на выбранных
бюджетах, а не побитовое совпадение всех промежуточных итератов.
Все масштабирования одной задачи являются зависимыми контролями.
Дополнительные исследования и полные настройки базовых методов
сохранены в расширенном техническом отчёте и сопроводительном архиве.

\begin{table}[htbp]
\centering\footnotesize
\setlength{\tabcolsep}{4pt}
\caption{Сертификация при $10^{-5}$ за 1000 обновлений. В ячейке указаны число успехов и медиана вызовов функции до остановки или лимита, включая неуспехи. Seoul: 30 стартов, Metro: 24 на радиус/метод. $F/A$: FW/AFW; индексы $1,U$: BT с начальной оценкой $1,U_0/D^2$; $Q$: короткий шаг с известными $H$. AC считает 1000 внешних проб, включая отказы, и оплачивает две предварительные исходные точки. Эти медианы не являются временем до успеха.}
\label{tab:counts}
\begin{tabular}{llrrrrrrr}
\toprule
Данные & $R/R_*$ & $F_1$ & $F_U$ & $A_1$ & $A_U$ & $F_Q$ & $A_Q$ & AC \\
\midrule
Seoul & 0.5 & 1; 2001 & 1; 2003 & 30; 424 & 30; 399 & 0; 1001 & 0; 1001 & 0; 1002 \\
Seoul & 2 & 0; 2001 & 0; 1996 & 0; 2001 & 0; 1995.5 & 0; 1001 & 0; 1001 & 0; 1002 \\
Metro & 0.5 & 3; 2001 & 3; 2003 & 24; 99 & 24; 80.5 & 1; 1001 & 24; 423 & 13; 385 \\
Metro & 2 & 0; 2000 & 0; 2002 & 6; 2000 & 7; 2001 & 0; 1001 & 0; 1001 & 0; 1002 \\
\bottomrule
\end{tabular}
\end{table}


\section{Что непосредственно устанавливают ближайшие оценки}
\label{app:comparison}
Следующее сравнение относится к конкретным сформулированным
результатам, а не ко всем возможным следствиям их доказательств.
Khademi и Silveti-Falls~\cite[теоремы 4.4--4.5]{khademi} сочетают
оценки по разности градиентов с обратным поиском. Их сформулированные
скорости для задач с ограничениями сохраняют вклад глобальной
гладкости/градиента; они также явно обсуждают медленную адаптацию
обычного обратного поиска к убывающей кривизне. Поэтому сама мотивация
известна ранее. Yagishita и Ito~\cite[алгоритм 3, теорема 3.1]{yagishita}
анализируют композитный условный градиент без линейного поиска с
неубывающей оценкой и невыпуклой гарантией для минимального зазора.
Это другие методы и утверждения об остановке, а не основание общей
претензии на приоритет. Pedregosa и соавторы~\cite[теорема 3]{pedregosa}
ограничивают прямой-двойственный зазор Фенхеля последнего итерата
через усреднённые принятые оценки кривизны. Их двойственная нижняя
величина и максимум скалярных касательных нижних границ не имеют
общего порядка для произвольных историй: у $f=x^2/2$ на $[-1,1]$
истории $(-1,1)$ и $(-1,-1/2)$ меняют порядок на противоположный.
Не утверждается, что эти истории являются траекториями алгоритма;
они не противоречат указанной теореме. Теорема 4 той работы уже
даёт геометрическое продвижение по принятой локальной кривизне,
которое специализирует наше AFW-доказательство.

Подстановка $L_Q=H_0+H_1\max_Q\delta$ в оценку глобальной
гладкости сохраняет глобальный диапазон зазора.
Меньший начальный сертификат $U_0$ не ограничивает кривизну на всём
$Q$; сужение области до начального множества подуровня изменило бы
её LMO. Наш аргумент первого возврата вместо этого ограничивает лишь
проверяемый отрезок спуска. Фазовый подсчёт контролирует оценки выше
текущего уровня.

Детерминированный выпуклый результат
\cite[приложение C.4]{lobanovkoloskova} уже разделяет режимы
зависящей от функции кривизны для градиентного спуска; это не теорема
об LMO-сертификате. Анализ FW при зависимости от нормы
градиента~\cite[теорема 5.1]{vyguzov} после прямой подстановки
содержит вклад максимального градиента. Ни использование фаз,
ни адаптация к неизвестной кривизне сами по себе не представлены
здесь как новые.

Takahashi и соавторы~\cite{bregmanfw} используют относительную
гладкость, ядро и масштабирование отрезков. Евклидово ядро с одной
глобальной $L_Q$ само по себе не даёт приведённого разделения.
В адаптивном слайдинге~\cite[теоремы 2.1 и 4.1]{adcgs} затраты
внутреннего LMO отличаются от внешних итераций, а теорема с
рестартами требует также локальной сильной выпуклости и параметров
расписания. Одно лишь условие~\eqref{eq:curvature} не даёт сильной
выпуклости: у $\cosh(x_1)-1$ на квадрате есть прямая минимизаторов.
Превосходство над слайдингом не установлено.

Wirth и соавторы~\cite[определение 2.3(b), теорема 3.6]{wirth}
используют равномерную расширенную кривизну и вершинную оценку
ошибки. Kerdreux и соавторы~\cite[утверждение 4.2, теорема 5.1]{restartfw}
уже сочетают начальную фазу, вычислимую остановку по сильному
зазору Вульфа, неизвестные параметры оценки ошибки и накопленные
вызовы LMO с глобальной кривизной шагов удаления.
Zhao~\cite[теоремы 4.1--4.2]{zhao} доказывает линейные скорости
по функции и FW-зазору для логарифмически однородных барьеров с
точным поиском или шагом по локальному гессиану. Наш класс функций
и проверяемый алгоритм отличаются. Ни одно из этих сравнений не
исключает эквивалентной уточнённой специализации другого адаптивного
доказательства.

\section{Воспроизводимость и проверки свидетельств}
\label{app:reproduction}
Сопроводительный пакет включает снимки исходников, атрибуцию сырых
данных, CSV-результаты и SHA-256-манифесты. Основное сравнение на
реальных данных объединяет 648 сохранённых записей со 108
дополнительными записями AC-FW на тех же условиях; они не считаются
новыми наборами данных. Исследование сертификатов повторяет 432
существующих пути и даёт 2160 зависимых причинных профилей.
Факторный эксперимент инициализации также повторно использует
прежние условия. Табличные значения генерируются из неизменяемых CSV;
реестр утверждений фиксирует каждую выборку и её значение.

Зафиксированная среда использует Python 3.13.5, NumPy 2.5.2,
SciPy 1.16.3, Matplotlib 3.10.9 и scikit-learn 1.8.0;
переменные среды числа потоков BLAS/OpenMP установлены в единицу.
Из корня извлечённого репозитория:
\begin{verbatim}
python -m pip install -r PART3/requirements.txt
python -m unittest discover -s PART3/Work -p "test_*.py"
python PART3/Work/validate_results.py
python PART3/Paper/prepare_article_evidence.py
\end{verbatim}
65 тестов проверяют производные, независимость от эталонов,
фактические вызовы оракулов, сертификаты, принятое продвижение и
пределы кривизны. Аудиты полных траекторий выполняются до
прореживания; сохранённые валидаторы проверяют также происхождение
и бюджетный учёт. Отдельный конечный аудит использует точную
рациональную арифметику для 8370 фазовых конфигураций, 57989
посещений абстрактных состояний кривизны, бюджетов фаз, обращения
зазора и подсчёта удаляемых атомов. Он охватывает надмножество
возможных историй кривизны и не утверждает реализуемость каждого
состояния целевой функцией. Конечные тесты не являются формальным
доказательством или независимой математической рецензией.

Семь дополнительных тестов проверяют точную дифференциальную
огибающую, её обращение и пределы больших параметров, одномерные
экстремизаторы в десятичной арифметике с 70 знаками и фактические
зазоры пяти прежних синтетических и регрессионных функций.
Отдельный рациональный архив содержит 24 AFW-запуска из равномерных
весов атомов на симплексах и прямоугольниках, шесть заданных локальных
состояний с 12 максимальными удалениями и 180 проверок квадратичных
проб положительного базового контрпримера. Все 30 запусков с нулевой
базовой кривизной достигают точно нулевых функции и сертификата.
Это конечные алгебраические проверки заявленных механизмов, а не
свидетельство универсального эмпирического ускорения. Оба архива
точного аудита повторно исполняются из собственного сохранённого кода.

Для симплексов $t=c^\top x$, $c=(0,0,1,2)$ и
$f=t+\beta t^3$, $\beta\in\{0,1/16,1\}$. Норма гессиана равна
$6\beta t\norm c^2=30\beta t\le30\beta f$ на симплексе.
Для прямоугольника $0\le t\le1$, $|y|\le1/2$ используется
$f=t+\beta t^3(1+y^2)$. Гессиан имеет неотрицательную диагональ и
определитель $12\beta^2t^4(1-2y^2)\ge0$; его норма не превосходит
следа $6\beta t(1+y^2)+2\beta t^3\le(19\beta/2)t\le10\beta f$.
Оба семейства имеют неодноточечную экспонированную оптимальную грань
и атомный отступ $m=1$. Замена $t^3$ на $(\max\{t,0\})^3$
даёт выпуклое $C^2$-продолжение на окрестность без изменения
допустимых значений. Четыре начальные оценки: $1/4,1,4,64$.
Заданные локальные состояния сохраняют два оптимальных и два
неоптимальных атома и завершаются двумя точными максимальными
шагами удаления. Ещё два рациональных пути в тестах начинаются
около этой грани с $M=65536$ и дают соответственно семь и восемь
немаксимальных шагов удаления перед последними двумя удалениями.
Они проверяют отдельный потенциал инициализации $Q_0$ в
теореме~\ref{thm:zerofinite}.

Более широкие исторические эксперименты содержат повторяющиеся
начальные вершины Bike Sharing: десять прежних seed дают семь
различных стартов. Поздние сводки перехода активности явно исключают
эти повторы. Исторические сводки с весами seed не используются как
независимые повторения в основном сравнении счётных данных, где
перебираются все различные вершины. Численные эталоны и аудиты
с плавающей точкой не дают гарантий направленного округления.

\paragraph{Раскрытие использования генеративных инструментов.}
Ассистент OpenAI Codex существенно использовался для первичного
отбора литературы, разработки доказательств, реализации, анализа
экспериментов, написания рукописи и автоматизации проверок.
Для критической проверки представлены явные выводы, снимки исходников
и исполняемые проверки. Автоматические проверки и внутренний анализ
с помощью ИИ не являются независимым математическим рецензированием.
Ответственность за научные утверждения подаваемой работы несут
указанные авторы.
